% !TEX root = ../../main.tex
% C2.1 - Existing commercial systems (translated). Every claim is sourced.
\section{Existing Systems on the Market}
\label{sec:2.1}

Three commercial products represent three approaches: comprehensive video analytics (BriefCam), appearance search (Avigilon) and natural-language search in the cloud (Verkada). Academic work is covered in Chapter~\ref{chapter:lythuyet}.

\subsection{BriefCam}
\label{sec:2.1.1}

BriefCam is video content analytics software based on deep learning, acquired by Canon in 2018~\cite{briefcam2018canon} and now distributed as a Milestone Systems product~\cite{milestone2026briefcam}. Its functions are split into three modules: \emph{Review} supports investigation on recorded video, filtering people and vehicles by attributes such as colour, size, speed and direction, finding objects that look similar to a sample object, and condensing hours of video into a short clip (Video Synopsis); \emph{Respond} raises real-time rule-based alerts; and \emph{Research} provides statistics, heat maps and people counting~\cite{milestone2026briefcam}. Figure~\ref{fig:briefcam} shows the search interface of the product. BriefCam can be installed on premises or in the cloud, but the processing server must have an NVIDIA graphics card~\cite{milestone2026briefcamfaq}.

\begin{figure}[H]
    \centering
    \includegraphics[width=0.6\textwidth]{briefcam.png}
    \caption{The search and filtering interface of BriefCam~\cite{milestone2026briefcam}}
    \label{fig:briefcam}
\end{figure}

\paragraph{Relation to the project.} BriefCam searches by attributes and by a sample object selected in recorded video, not an uploaded image as in this project; its documentation does not mention natural-language search.

\subsection{Avigilon Appearance Search}
\label{sec:2.1.2}

Avigilon Appearance Search is the deep-learning video search technology of Avigilon, part of Motorola Solutions, used to find a person or vehicle across hours of video from many cameras~\cite{motorola2021appearance}. A search can start by choosing appearance features from a predefined list such as clothing colour, gender and age group; by uploading an image; or by selecting an object in the video as a sample~\cite{motorola2021appearance}. The results help trace the route of the subject and can be bookmarked and exported as an evidence package describing the event. Figure~\ref{fig:avigilon} shows the interface with a person marked. The technology is built into the Avigilon Control Center video management software and runs on the vendor's cameras and recorders, or on the Avigilon AI Appliance for third-party ONVIF cameras.

\begin{figure}[H]
    \centering
    \includegraphics[width=0.45\textwidth]{avigilon.png}
    \caption{Avigilon Control Center with a marked subject~\cite{avigilon2026appearance}}
    \label{fig:avigilon}
\end{figure}

\paragraph{Relation to the project.} Avigilon supports two of the project's forms of description, attributes and a reference image, and compiles results into evidence like the project's case files; free-text search is not mentioned, and it is tied to the vendor's software and devices.

\subsection{Verkada AI-Powered Search}
\label{sec:2.1.3}

Verkada AI-Powered Search, introduced in 2024 on the Command cloud security platform, lets users search for people and vehicles in everyday language, for example ``person wearing helmet'', instead of choosing from a list of attributes~\cite{verkada2024aisearch}, and search with an uploaded image of a person, vehicle or object~\cite{verkada2024reverse}. According to the vendor's technical article~\cite{verkada2025scaling}, cameras produce crops of objects; in the cloud, the crops are encoded into vectors with the CLIP model, which maps images and text into one representation space, and stored in a vector database; the query sentence is encoded the same way to find the nearest crops. Figure~\ref{fig:verkada} shows the results of a natural-language query. The system works best with Verkada cameras; third-party cameras connect through a Command Connector device with higher latency~\cite{verkada2024connector}.

\begin{figure}[H]
    \centering
    \includegraphics[width=0.52\textwidth]{verkada.png}
    \caption{Results of a natural-language query on the Verkada Command platform~\cite{verkada2024aisearch}}
    \label{fig:verkada}
\end{figure}

\paragraph{Relation to the project.} Verkada is technically closest: both search by description and by image through a shared representation space and a vector database. Verkada, however, runs on the vendor's cloud, while the project runs on premises on a commodity computer.
